
Low-Rank Adaptation (LoRA) fine-tunes large language models efficiently but typically assigns uniform rank across all layers, ignoring heterogeneous computational roles that emerge during pre-training. Existing adaptive rank methods learn the allocation during training, introducing a circular dependency and per-experiment overhead. This work asks whether a pre-trained model already encodes a reliable signal for rank allocation decisions prior to any fine-tuning.

The layer-wise MLP activation sparsity profile of LLaMA-3.1-8B is characterized as a structural fingerprint measurable from a single forward pass. Using forward hooks on 512 calibration samples, the fingerprint exhibits significant heterogeneity across 32 layers (CV = 0.544 at epsilon = 0.01) and is stable across four diverse calibration distributions: intraclass correlation ICC(3,k) = 0.9846 (95\% CI: [0.97, 0.99]), with all six pairwise Kendall's tau values at or above 0.734. Threshold invariance is also confirmed: all six cross-epsilon Kendall's tau values exceed 0.95 across epsilon values spanning two orders of magnitude (minimum observed: tau = 0.9597 for the epsilon = 0.01 vs. epsilon = 0.1 pair).

These results establish that LLaMA-3.1-8B's sparsity profile reflects pre-training geometry rather than input content, providing an empirical foundation for its use as a zero-cost structural prior for LoRA rank allocation. Whether this fingerprint predicts rank requirements and enables efficient allocation remains empirically open; the corresponding experiments (H-M3, H-M4) were not executed in this work.

